% !PRINTED 306
\weberpage{306}
%<<CONT
\[
\begin{aligned}
S_3 &= S_4 + \begin{pmatrix} -1, & 1 \\ 1, & 1 \end{pmatrix} S_4 \\
S_2 &= S_3 + \begin{pmatrix} 0, & 1 \\ -1, & 0 \end{pmatrix} S_3 \\
S_1 &= S_2 + \begin{pmatrix} 1, & 0 \\ 1, & 1 \end{pmatrix} S_2 + \begin{pmatrix} 1, & 0 \\ -1, & 1 \end{pmatrix} S_2 \\
S &= S_1 + \begin{pmatrix} -1, & 0 \\ 0, & 1 \end{pmatrix} S_1.
\end{aligned}
\]

Wir haben hiernach die Compositions- und Indexreihe von $S$

\[
\begin{gathered}
S,\ S_1,\ S_2,\ S_3,\ S_4, \\
2 \quad 3 \quad 2 \quad 2
\end{gathered}
\]

\gesperrt{und wir kommen also hier zu dem Ergebniss, dass alle Gleichungen
 9\textsuperscript{ten} Grades mit linearer Congruenzgruppe
metacyklisch sind.}

\begin{center}
\S.\,79.

\gesperrt{Tern\"are lineare Congruenzgruppe f\"ur den Modul 2.}
\end{center}

Mannigfache interessante Beziehungen ergeben sich, wenn
wir die tern\"are lineare Congruenzgruppe $R$ f\"ur den Modul 2
betrachten, die als transitive Permutationsgruppe von acht Elementen
 aufgefasst werden kann, und die Gruppe der metacyklischen
 Gleichungen 8\textsuperscript{ten} Grades enthalten muss. Nach dem
oben Bewiesenen kommt es vor Allem darauf an, die homogene
Congruenzgruppe $S$ zu betrachten, die aus den Elementen

\weq{1}{\sigma = \begin{pmatrix} a, & b, & c \\ a_1, & b_1, & c_1 \\ a_2, & b_2, & c_2 \end{pmatrix}}

besteht, worin die Ziffern $a, b, c, \ldots$ nach dem Modul 2, also $= 0$
oder $= 1$, anzunehmen sind, und die nach der im \S.\,37 gegebenen
Vorschrift componirt werden. Es kommen dabei nur solche
Systeme der Zahlen $a, b, c, \ldots$ in Betracht, deren Determinante
$= 1$, d.\,h. ungerade ist.

Um den Grad der Gruppe $S$ zu ermitteln, beachte man, dass
die erste Zeile von $\sigma$ sieben verschiedene Formen haben kann:

\[
(1,0,0),\ (0,1,0),\ (0,0,1),\ (0,1,1),\ (1,0,1),\ (1,1,0),\ (1,1,1).
\]
